% Options for packages loaded elsewhere
\PassOptionsToPackage{unicode}{hyperref}
\PassOptionsToPackage{hyphens}{url}
\PassOptionsToPackage{dvipsnames,svgnames,x11names}{xcolor}
%
\documentclass[
  10pt,
]{article}
\usepackage{amsmath,amssymb}
\usepackage{iftex}
\ifPDFTeX
  \usepackage[T1]{fontenc}
  \usepackage[utf8]{inputenc}
  \usepackage{textcomp} % provide euro and other symbols
\else % if luatex or xetex
  \usepackage{unicode-math} % this also loads fontspec
  \defaultfontfeatures{Scale=MatchLowercase}
  \defaultfontfeatures[\rmfamily]{Ligatures=TeX,Scale=1}
\fi
\usepackage{lmodern}
\ifPDFTeX\else
  % xetex/luatex font selection
  \setmainfont[]{Noto Serif}
  \setmonofont[]{DejaVu Sans Mono}
\fi
% Use upquote if available, for straight quotes in verbatim environments
\IfFileExists{upquote.sty}{\usepackage{upquote}}{}
\IfFileExists{microtype.sty}{% use microtype if available
  \usepackage[]{microtype}
  \UseMicrotypeSet[protrusion]{basicmath} % disable protrusion for tt fonts
}{}
\makeatletter
\@ifundefined{KOMAClassName}{% if non-KOMA class
  \IfFileExists{parskip.sty}{%
    \usepackage{parskip}
  }{% else
    \setlength{\parindent}{0pt}
    \setlength{\parskip}{6pt plus 2pt minus 1pt}}
}{% if KOMA class
  \KOMAoptions{parskip=half}}
\makeatother
\usepackage{xcolor}
\usepackage[margin=24mm]{geometry}
\setlength{\emergencystretch}{3em} % prevent overfull lines
\providecommand{\tightlist}{%
  \setlength{\itemsep}{0pt}\setlength{\parskip}{0pt}}
\setcounter{secnumdepth}{-\maxdimen} % remove section numbering
\ifLuaTeX
  \usepackage{selnolig}  % disable illegal ligatures
\fi
\usepackage{bookmark}
\IfFileExists{xurl.sty}{\usepackage{xurl}}{} % add URL line breaks if available
\urlstyle{same}
\hypersetup{
  pdftitle={Nonsingularity of a Logarithmic Coefficient Matrix for Centrally Symmetric Seven-Node Rational Configurations},
  pdfauthor={Ryutaro Yonezu},
  colorlinks=true,
  linkcolor={Maroon},
  filecolor={Maroon},
  citecolor={Blue},
  urlcolor={Blue},
  pdfcreator={LaTeX via pandoc}}

\title{Nonsingularity of a Logarithmic Coefficient Matrix for Centrally
Symmetric Seven-Node Rational Configurations}
\usepackage{etoolbox}
\makeatletter
\providecommand{\subtitle}[1]{% add subtitle to \maketitle
  \apptocmd{\@title}{\par {\large #1 \par}}{}{}
}
\makeatother
\subtitle{Research No.~81 --- AI-audited proof-source preprint
candidate}
\author{Ryutaro Yonezu}
\date{2026-11-22}

\begin{document}
\maketitle

{
\hypersetup{linkcolor=}
\setcounter{tocdepth}{3}
\tableofcontents
}
\textbf{Research status:} General theorem: \textbf{independent AI
mathematics audit PASS}, as reported in the Stage 06 addendum;
six-configuration Windows exact finite check: \textbf{PASS (not a full
theorem audit)}. \textbf{NOVELTY INCONCLUSIVE / REFERENCES PARTIAL / NOT
HUMAN PEER REVIEWED / All rights reserved.}

\textbf{Editorial notice:} This is a collected proof-source edition. The
five mathematical source documents below are reproduced verbatim, with
no change to their mathematical statements, equations or proofs.
Statements in the original dated source documents such as ``PRIVATE /
HOLD'', ``pending audit'', and ``no Windows test'' record their
\emph{earlier drafting status}. The current audit status is given above;
it does not imply human peer review or a separate mathematical audit of
this typesetting.

\textbf{Provenance:} Frozen v0.7 SHA256
\texttt{ef40748e44e13280e00b4db496a3296119fa1adb2659afb869e5865f3765c6f4}.
The original six-stage outer parent ZIP SHA256
\texttt{2cfcad697ae600b1f3f4ca9b925229d999d0384c0d3ea1c78d9b7fab50fe06ee}
was recovered and directly SHA256-verified; the Stage 05 audit request
ZIP SHA256
\texttt{a4317074478a5af590eda36bc66c3aee682e0ee64f013024d5805e52ea678419}
was also recovered and verified. The original completion proof and
certificate match those source archives byte-for-byte. The finite
Windows check does not prove the general theorem. The optional
single-GCD source is supporting historical material and is not needed in
the final proof chain. See the accompanying reproducibility materials
for the D=2 reduction and the supporting certificates.

\section{Main theorem and final completion
argument}\label{main-theorem-and-final-completion-argument}

\section{Seven-node rational regularity via a seventh-power arithmetic
progression}\label{seven-node-rational-regularity-via-a-seventh-power-arithmetic-progression}

2026-11-20 JST --- Research continuation, PRIVATE / HOLD. NOVELTY
INCONCLUSIVE / REFERENCES PARTIAL / NOT HUMAN PEER REVIEWED / All rights
reserved.

\subsection{Statement and
dependencies}\label{statement-and-dependencies}

The original D=2 logarithmic coefficient matrix A, with entries
A{[}n,(i,j){]} = b\_i\^{}(n-j)/(n-j) for n\textgreater j and zero
otherwise (n=1,\ldots,18; i=1,\ldots,6; j=0,1,2), is nonsingular for the
seven distinct nodes \{0, ±a, ±b, ±c\}, where a,b,c are positive
distinct rationals. This is a proof proposal pending independent audit.
Scaling by a nonzero rational lambda multiplies row n by lambda\^{}n and
column (i,j) by lambda\^{}(-j), so denominators may be cleared and the
gcd removed. Order 0\textless a\textless b\textless c and gcd(a,b,c)=1.

Dependencies: frozen v0.7 (SHA256
ef40748e44e13280e00b4db496a3296119fa1adb2659afb869e5865f3765c6f4), in
particular its nested v0.5 necessary residue reduction;
symmetric\_D2\_radius\_proof.md; symmetric\_seven\_2adic\_proof.md;
symmetric\_seven\_small\_prime\_proof.md; Theorem 1 in
symmetric\_seven\_global\_proof.md. No reverse residue reduction is
used. The preceding single-gcd case analysis is not needed. The cited
theorem of Darmon and Merel is an external established theorem, not
newly proved here.

\subsection{Proof}\label{proof}

Suppose A singular. By the global factorization theorem, define
d=gcd(b,c), e=gcd(a,c), f=gcd(a,b). Then there exist positive integers
u,v,w,g and eps\_a,eps\_b,eps\_c in \{1,2\}, exactly one 2, with

\begin{verbatim}
a=e f u^2, b=d f v^2, c=d e w^2,
c^2-b^2=g eps_a d^2 u^3,
c^2-a^2=g eps_b e^2 v^3,
b^2-a^2=g eps_c f^2 w^3.
\end{verbatim}

The three square factors u,v,w have common gcd 1: any common prime would
divide a,b,c.~After dividing each difference by its shared factor
squared,

\begin{verbatim}
e^2 w^4-f^2 v^4=g eps_a u^3,
d^2 w^4-f^2 u^4=g eps_b v^3,
d^2 v^4-e^2 u^4=g eps_c w^3.
\end{verbatim}

Multiply the first by u\^{}4, the second by v\^{}4, the third by w\^{}4.
The sum of the first and third left sides equals the second. Cancel
positive g to obtain

\begin{verbatim}
eps_a u^7 + eps_c w^7 = eps_b v^7.
\end{verbatim}

Darmon--Merel, Main Theorem (1), states that X\textsuperscript{n+Y}n=2
Z\^{}n has no nontrivial primitive integer solution for n\textgreater=3,
where trivial means XYZ=0 or ±1. The author-hosted manuscript dated
2001-01-05, pp.~1--2, was consulted directly (including page images);
the journal article is J. reine angew. Math. 490 (1997), 81--100, DOI:
10.1515/crll.1997.490.81. URL:
https://perso.imj-prg.fr/loic-merel/wp-content/uploads/merel-pub/winding.pdf
.

If eps\_a=2 then (X,Y,Z)=(v,-w,u); if eps\_c=2 then (v,-u,w); if
eps\_b=2 then (u,w,v). Each triple is primitive and has nonzero product.
The theorem forces u=v=w=1, because they are positive integers. The
first two sign-mixed cases are then incompatible with the original
seventh-power equality. Hence eps\_b=2, eps\_a=eps\_c=1.

The divided square-difference equations now reduce to
e\textsuperscript{2-f}2=g and d\textsuperscript{2-e}2=g. Put
t=f\^{}2\textgreater0 and h=g\textgreater0. Then the squared radii
(X,Y,Z)=(a\textsuperscript{2,b}2,c\^{}2) are

\begin{verbatim}
X=t(t+h), Y=t(t+2h), Z=(t+h)(t+2h).
\end{verbatim}

No assumption about which shared factors are nontrivial is needed. In
particular this remaining branch must be checked, rather than being
dismissed as a trivial configuration: distinct radii are still possible
before imposing the scalar equations.

For squared radii X\_i in \{X,Y,Z\}, put

\begin{verbatim}
g0=1/(-XYZ)^3,
gi=1/[8 X_i^3 product_{j!=i}(X_i-X_j)^3],
G=g0+2 sum_i gi,
C=-sum_i gi (9+12 X_i sum_{j!=i}1/(X_i-X_j)).
\end{verbatim}

The centered D=2 residue matrix M has det M=-(G+C)(3G+2C)/8. Set
H=h\^{}6 t\textsuperscript{6(t+h)}6(t+2h)\^{}6. Exact rational
simplification yields

\begin{verbatim}
G+C=3 F/(32 H),
3G+2C=3 J/(4 H),
\end{verbatim}

where

\begin{verbatim}
F=84h^6+600h^5t+1396h^4t^2+1680h^3t^3+1150h^2t^4+438ht^5+73t^6,
J=(11h^2+12ht+6t^2)(2h^4+12h^3t+18h^2t^2+12ht^3+3t^4).
\end{verbatim}

All factors and coefficients are positive at t,h\textgreater0. Both
rational identities are obtained by direct substitution into the scalar
expressions and checked symbolically by
symmetric\_seven\_completion\_certificate.py. Thus det M\textless0 on
this branch, contradicting singularity of A, which requires singularity
of M. QED subject to the stated proof dependencies.

\subsection{Status and limits}\label{status-and-limits}

This closes all shared-gcd patterns simultaneously, including the
arithmetic-progression exceptional branch. It is an unbounded theorem
proposal for the centrally symmetric seven-node D=2 family only. It does
not settle arbitrary nonsymmetric seven-node configurations, nine-node
symmetric configurations, arbitrary D or the original unrestricted
rational conjecture. The new proof uses exact symbolic identities, no
finite height cutoff, and depends on the external generalized Fermat
theorem. It is not an independent audit. The frozen mathematical source
is unchanged; no v0.8 freeze or Windows test; no publication permission.
Novelty remains inconclusive. The primewise support proof and necessary
residue reduction require review as part of the future independent
audit.

\newpage

\section{Supporting source 1:
symmetric\_seven\_2adic\_proof.md}\label{supporting-source-1-symmetric_seven_2adic_proof.md}

\section{A 2-adic certificate for centrally symmetric seven-node
configurations
(D=2)}\label{a-2-adic-certificate-for-centrally-symmetric-seven-node-configurations-d2}

PRIVATE / HOLD --- NOVELTY INCONCLUSIVE / REFERENCES PARTIAL / NOT HUMAN
PEER REVIEWED / All rights reserved.

Dependencies: frozen v0.7 and \texttt{symmetric\_D2\_radius\_proof.md}
(necessary residue reduction and radius formulas). No inverse residue
reduction is invoked. Normalize rational radii to primitive distinct
positive integers a,b,c by scaling. This note is a computer-assisted
proof proposal, not independently audited.

\subsection{Theorem}\label{theorem}

The original matrix is nonsingular if at least one primitive radius is
even, or if all three are odd and the sorted 2-adic valuations of their
squared differences are d,d,e with e\textgreater=d+2. Thus any singular
rational seven-node centrally symmetric configuration must have three
odd primitive radii with depths d,d,d+1, d\textgreater=3.

\subsection{Polynomial proof}\label{polynomial-proof}

Use X,Y,Z for the squared radii. With g0=(-XYZ)\^{}(-3), gi=1/(8 Xi\^{}3
product\_\{j!=i\}(Xi-Xj)\^{}3), set G=g0+2 sum gi and C=-sum gi(9+12 Xi
sum\_\{j!=i\}1/(Xi-Xj)). The previous reduction gives det
M=-(G+C)(3G+2C)/8. Construct primitive integer numerator polynomials
P1,P2 of G+C,3G+2C by exact cancellation over QQ, selecting signs so the
coefficients of A\^{}4 in Pi(1,A,0) are -5,-3 respectively. Each
polynomial is symmetric, homogeneous of degree 12, and has 61 nonzero
monomials. Denominators are nonzero for the allowed squared radii.
Homogeneity allows division by any odd squared radius, which is a 2-adic
unit.

For a term c\_ij A\^{}i B\^{}j write W=i e+j f+v2(c\_ij), where e=v2(A),
f=v2(B). The companion script constructs the polynomials from the
rational functions and verifies all stated strict weighted inequalities
for all monomials. The verification is finite because W is affine in
e,f; each unbounded cone is given as a base plus nonnegative
combinations of generators, and strict positivity at the base plus
nonnegative slopes suffices. It does not search just a bounded set of
valuations.

\begin{itemize}
\tightlist
\item
  One odd radius, two even with unequal valuations: use (1,A,B),
  e\textgreater=2, f\textgreater=e+2. The unique lowest term is A\^{}4
  for both polynomials; valuations are 4e. Cone base (2,4), generators
  (1,1),(0,1).
\item
  One odd radius, two even with equal valuations: use (1,A,A+B),
  e\textgreater=2, f\textgreater=e+3. The latter inequality follows
  since odd squares are 1 mod 8 after removing equal powers of 2. A\^{}4
  is uniquely lowest, with coefficient 8 or 4 respectively; valuations
  4e+3 and 4e+2. Cone base (2,5), generators (1,1),(0,1).
\item
  Two odd radii and one even: use (1,1+A,B), e\textgreater=3,
  f\textgreater=2 even. Only candidates for the least W are A\^{}4 and
  B\^{}3, with respective coefficients (-5,-8) for P1 and (-3,-4) for
  P2. Every other term has larger W than at least one candidate, checked
  at base (3,2) with coordinatewise nonnegative slopes. P1 has unequal
  candidate valuations 4e and 3f+3 by parity. For P2, the tie 4e=3f+2
  occurs exactly at (e,f)=(5+3t,6+4t), t\textgreater=0. All other terms
  have W at least 4e+3, verified at t=0 with nonnegative slopes. For
  A=2\^{}e u, B=2\^{}f v\^{}2 with u,v odd 2-adic units, the tied terms
  after division by 2\^{}(4e) give -3u\textsuperscript{4-v}6 = -4 mod 8
  (odd square gives v\^{}6=1 mod 8). So the sum has valuation exactly
  4e+2, without vanishing.
\item
  All odd radii: by the ultrametric inequality and the two-element
  residue field the three squared-difference depths are d,d,e with
  e\textgreater d, d\textgreater=3. Permute so the closest pair is
  (1,1+A), and use (1,1+A,1+B) with v2(A)=e,v2(B)=d.~If
  e\textgreater=d+2, the unique lowest term is B\^{}5, coefficient 8 for
  P1, 4 for P2. Cone base (5,3), generators (1,1),(1,0). Hence both
  polynomials are nonzero. This completes the claimed sufficient
  conditions.
\end{itemize}

\subsection{A genuine limitation of the
method}\label{a-genuine-limitation-of-the-method}

Put (X,Y,Z)=(1,1+16u,9). The polynomial P1 divided by 2\^{}18 and P2
divided by 2\^{}17 are integer polynomials in u, each reducing to u+1
mod 2. By simple-root Hensel lifting each has its own root u in Z\_2
with u odd. Then 1+16u is a nonzero 2-adic square (every unit congruent
to 1 mod 8 is a square in Q\_2). The squared radii are distinct since
their difference depths are 3,3,4. Thus each scalar individually admits
a genuine Q\_2 zero in the unresolved regime; the roots need not
coincide. Consequently the reduced determinant can vanish over Q\_2.
This alone does not prove singularity of the original matrix over Q\_2,
since the converse reduction was not invoked, and is not a rational
counterexample.

\subsection{Finite checks and
limitations}\label{finite-checks-and-limitations}

Run \texttt{symmetric\_seven\_2adic\_certificate.py} beside
\texttt{symmetric\_D2\_radius\_certificate.py},
\texttt{d2\_moment\_reduction.py} and the frozen v0.7 ZIP. Requires
Python/SymPy. All weighted inequalities and both Hensel polynomial
reductions verified; lifts checked mod 2\^{}17. Primitive triples with
maximum radius \textless=50: 16648 total; new theorem covers 15418 (all
checked by exact rational scalar evaluation); old radius witnesses for
primes 5\textless=p\textless101 cover 16628; union 16643, of which 15
are newly covered. Remaining five verified by both frozen exact M and
original integer-cleared determinant. Frozen SHA256 verified. No
exhaustive rerun of previous certificates, no fresh literature search,
no independent audit or Windows validation. No assertion for arbitrary
radii in the unresolved regime or for arbitrary nonsymmetric seven-node
configurations.

\newpage

\section{Supporting source 2:
symmetric\_seven\_small\_prime\_proof.md}\label{supporting-source-2-symmetric_seven_small_prime_proof.md}

\section{Small-prime nonvanishing for centrally symmetric seven-node
configurations
(D=2)}\label{small-prime-nonvanishing-for-centrally-symmetric-seven-node-configurations-d2}

PRIVATE / HOLD --- NOVELTY INCONCLUSIVE / REFERENCES PARTIAL / NOT HUMAN
PEER REVIEWED / All rights reserved.

\subsection{Scope and dependencies}\label{scope-and-dependencies}

This is a computer-assisted proof proposal extending
\texttt{symmetric\_D2\_radius\_proof.md} and
\texttt{symmetric\_seven\_2adic\_proof.md}. The frozen v0.7 necessary
residue reduction is assumed; its converse is not needed. Normalize
rational radii to distinct positive primitive integers a,b,c.~Neither
the frozen bundle nor previous certificates are modified. No independent
audit or Windows validation.

\subsection{Theorem}\label{theorem-1}

For p=3 or p=5 the original D=2 matrix for nodes \{0,±a,±b,±c\} is
nonsingular if at least one radius is divisible by p, or if the three
p-adic valuations of squared-radius differences are not all equal.

Thus, using also the earlier 2-adic result, any singular rational
configuration of this type must have all three primitive radii coprime
to 30, squared difference depths d,d,d+1 at 2 (d\textgreater=3), and all
three squared-difference depths equal at 3 and at 5. At 3, the common
depth is \textgreater=1 because all nonzero squares are 1 modulo 3. At 5
the common depth is also \textgreater=1: three unit squares occupy at
most two residue classes. These are necessary conditions, not a full
nonsingularity theorem.

\subsection{Polynomial proof}\label{polynomial-proof-1}

Let X,Y,Z denote squared radii. Define g0=(-XYZ)\^{}(-3), gi=(8 Xi\^{}3
product\_\{j!=i\}(Xi-Xj)\textsuperscript{3)}(-1), G=g0+2 sum gi, C=-sum
gi(9+12 Xi sum\_\{j!=i\}1/(Xi-Xj)). The earlier radius identity gives
det M=-(G+C)(3G+2C)/8. Take the primitive integer numerators P1,P2 of
these two scalars, with signs fixed by the coefficients -5,-3 of A\^{}4
in Pi(1,A,0). Both are symmetric homogeneous polynomials of degree 12
with 61 terms; denominators are nonzero on admissible radii. By
homogeneity and symmetry normalize a unit squared radius to 1. All other
coordinates remain p-adic squares. The companion certificate verifies
strict weighted monomial comparisons on unbounded cones: if
W(i,j)=ie+jf+v\_p(c\_ij), strict inequality at a base and nonnegative
slope along each cone generator imply the inequality throughout the
cone. The polynomial identities and comparisons below are reconstructed
from the rational scalar expressions, not from a hard-coded list of
coefficients.

\subsubsection{Two divisible radii}\label{two-divisible-radii}

Use (1,A,B), e=v\_p(A), f=v\_p(B), both even and \textgreater=2. If
unequal, arrange f\textgreater=e+2. The unique least-weight term is
A\^{}4 with coefficient -5 or -3 respectively (including its coefficient
valuation). The certificate checks base (2,4) and rays (1,1),(0,1).

If e=f, divide A,B by p\^{}e to get square units u,v. The homogeneous
degree-four parts are

H1=-(A\textsuperscript{2-3AB+B}2)(5A\textsuperscript{2-2AB+5B}2),
H2=-(A\textsuperscript{2-3AB+B}2)(3A\textsuperscript{2-2AB+3B}2).

Every other term has degree \textgreater=5, hence valuation
\textgreater=5e. For p=3 the unit squares both reduce to 1 and both
H\_i(u,v) are units. For p=5 and u=v mod 5 both are units. For p=5 and
u!=v mod 5 the residues are 1,4 in some order. The second factor of each
H\_i is a unit and the first factor has valuation exactly one by

4(u\textsuperscript{2-3uv+v}2)=(2u-3v)\textsuperscript{2-5v}2,

since 2u-3v is divisible by 5. Thus the initial valuation is 4e or 4e+1,
strictly less than 5e. Neither scalar numerator vanishes.

\subsubsection{One divisible radius}\label{one-divisible-radius}

If the two unit squares agree modulo p, use (1,1+A,B),
e=v\_p(A)\textgreater=1, f=v\_p(B)\textgreater=2 even. The only possible
lowest terms are -5 A\^{}4 and -8 B\^{}3 for P1, or -3 A\^{}4 and -4
B\^{}3 for P2. The certificate compares each other term against at least
one of these two on the cone e\textgreater=1,f\textgreater=2. If the
coefficient of A\^{}4 has valuation 1, equality of weights is
impossible: 4e+1 is odd and 3f even. Otherwise, in case of a tie, the
unit part after scaling equals -5u\textsuperscript{4-8v}6 or
-3u\textsuperscript{4-4v}6, where u is an arbitrary unit and v is a unit
square root of the unit part of B. Direct enumeration of the nonzero
residues for p=3,5 in the cases with unit A\^{}4 coefficient proves
nonvanishing. If the two unit squares differ modulo p, necessarily p=5.
Use (1,4+A,B). The constant coefficients -103680 and -62208 have
valuations 1 and 0; the certificate verifies they are uniquely lowest
for e\textgreater=1,f\textgreater=2.

\subsubsection{No divisible radii}\label{no-divisible-radii}

For p=5 if two square classes occur, exactly one pair collides; the
earlier mirror collision theorem applies (p\textgreater3). Otherwise all
squares occupy one residue class (always the case at p=3). If pairwise
difference depths are unequal, the ultrametric property gives d,d,e with
e\textgreater d\textgreater=1. Choose the closest pair as (1,1+A), and
the remaining square as 1+B, with v\_p(A)=e and v\_p(B)=d.~The unique
least-weight term in Pi(1,1+A,1+B) is B\^{}5, with coefficient 8 or 4.
The certificate verifies base (2,1) and rays (1,1),(1,0). Both
coefficients are units for p=3,5. This completes the proof given the
earlier necessary residue reduction.

\subsection{Verification and
limitations}\label{verification-and-limitations}

Run \texttt{symmetric\_seven\_small\_prime\_certificate.py} in the
directory containing \texttt{symmetric\_D2\_radius\_certificate.py},
\texttt{d2\_moment\_reduction.py} and the frozen v0.7 ZIP. Requires
Python/SymPy. Both scalar polynomials passed all unbounded weighted cone
and residue checks. For primitive positive triples with maximum radius
\textless=50: 16648 total, p=3 covers 15236, p=5 covers 15526, their
union covers 16546; all 16546 checked against exact rational scalar
evaluations. The previous 2-adic and odd-prime witness conditions
(primes 5\textless=p\textless101) cover 16643; all conditions together
cover 16648. Previous five uncovered triples all covered by the new
conditions. The original determinant was not re-evaluated for this full
sample, nor was the previous weighted certificate rerun: old witness
coverage is recalculated using its predicate. Frozen SHA256 verified. No
novel claim about the unresolved primitive triples in general, nor about
nonsymmetric nodes. No new literature search or independent audit.

\newpage

\section{Supporting source 3:
symmetric\_seven\_global\_proof.md}\label{supporting-source-3-symmetric_seven_global_proof.md}

\section{Global factorization obstruction for symmetric seven-node
rational configurations
(D=2)}\label{global-factorization-obstruction-for-symmetric-seven-node-rational-configurations-d2}

PRIVATE / HOLD --- NOVELTY INCONCLUSIVE / REFERENCES PARTIAL / NOT HUMAN
PEER REVIEWED / All rights reserved.

\subsection{Scope and dependencies}\label{scope-and-dependencies-1}

Use the frozen v0.7 necessary residue reduction and the symmetric D=2
scalar formulas in \texttt{symmetric\_D2\_radius\_proof.md}. Use the
prior 2-adic and p=3,5 theorems in
\texttt{symmetric\_seven\_2adic\_proof.md} and
\texttt{symmetric\_seven\_small\_prime\_proof.md}. Both corresponding
certificates were replayed successfully this turn. The reduction's
converse is not required. The frozen v0.7 ZIP is unmodified, SHA256
ef40748e44e13280e00b4db496a3296119fa1adb2659afb869e5865f3765c6f4. No
independent audit or Windows test.

\subsection{Theorem 1 (global necessary
factorization)}\label{theorem-1-global-necessary-factorization}

Let 0\textless a\textless b\textless c be primitive positive integer
radii, obtained by clearing rational denominators and removing their
common gcd. Suppose the original D=2 matrix at \{0,±a,±b,±c\} is
singular. Set d\_a=gcd(b,c), d\_b=gcd(a,c), d\_c=gcd(a,b). Then there
are positive integers u\_a,u\_b,u\_c such that

\begin{verbatim}
a=d_b d_c u_a^2, b=d_a d_c u_b^2, c=d_a d_b u_c^2.
\end{verbatim}

Let
g=gcd(c\textsuperscript{2-b}2,c\textsuperscript{2-a}2,b\textsuperscript{2-a}2).
There are epsilon\_a,epsilon\_b,epsilon\_c in \{1,2\}, exactly one equal
to 2, such that

\begin{verbatim}
c^2-b^2 = g epsilon_a d_a^2 u_a^3,
c^2-a^2 = g epsilon_b d_b^2 u_b^3,
b^2-a^2 = g epsilon_c d_c^2 u_c^3.
\end{verbatim}

Furthermore gcd(g,abc)=1, and all radii are coprime to 30. These
identities are necessary, not sufficient, for singularity.

\subsubsection{Proof}\label{proof-1}

Let P\_1,P\_2 be the signed primitive numerators of G+C, 3G+2C from the
preceding certificates. They are symmetric homogeneous integer
polynomials in the squared radii. Singular A implies singular M, hence
at least one P is zero. Fix such a P. Set (K,J)=(5,8) or (3,4),
respectively. Direct symbolic reconstruction in the companion
certificate proves the following identities/support facts:

\begin{enumerate}
\def\labelenumi{\arabic{enumi}.}
\tightlist
\item
  P(1,Y,0)=-K Y\textsuperscript{4(Y-1)}4; P(1,1,Z)=J
  Z\textsuperscript{3(Z-1)}5.
\item
  P(1,1+A,B) has precisely the componentwise-minimal monomials -K A\^{}4
  and -J B\^{}3. Every other monomial is strictly heavier than one of
  them for positive valuations of A,B at p\textgreater=7.
\item
  P(1,A,B) has no total degree below four, and its degree-four part is
  -(A\textsuperscript{2-3AB+B}2)(K A\^{}2-2AB+K B\^{}2). At A=B=1 this
  part equals J. For unequal positive valuations of A,B, its
  least-weight term is the pure fourth power of the less-divisible
  variable, strictly below every other term, for p\textgreater=7.
\item
  In P(1,1+A,1+B), for e=v\_p(A)\textgreater=f+1 where
  f=v\_p(B)\textgreater=1, the term J B\^{}5 is uniquely least-weight at
  all odd p for which J is a unit. Indeed every other exponent (i,j)
  satisfies i+j\textgreater=5 and 2i+j\textgreater5; hence
  ie+jf\textgreater5f.
\end{enumerate}

These follow from the finite list of monomial exponents, with identities
verified exactly in \texttt{symmetric\_seven\_global\_certificate.py}.
In each p-adic chart divide squared radii by a unit squared radius,
using symmetry and homogeneity; none of the valuation comparisons
changes when multiplying by such units. The normalized coordinates are
p-adic integers.

At p=2 the preceding theorem forces odd radii and sorted difference
depths d,d,d+1. At p=3,5 the preceding theorems force unit radii and
equal difference depths. Thus these primes never divide a radius.

For any p\textgreater=7:

\begin{itemize}
\tightlist
\item
  If p divides no radius, the three difference depths must be equal. If
  just two squares collide modulo p, identity 1 on P(1,1,Z) excludes
  zero. If all three collide with unequal depths, choose the closest
  pair as (1,1+A) and the third square as 1+B; fact 4 excludes zero. If
  no pair collides all depths are zero.
\item
  If p divides exactly one radius, say a, singularity forces
  b\textsuperscript{2=c}2 mod p by identity 1 applied to P(1,Y,0).
  Choose b\^{}2 as the unit normalizer. In P(1,1+A,B),
  v\_p(A)=v\_p(c\textsuperscript{2-b}2)=e\textgreater0 and
  v\_p(B)=2v\_p(a)=f\textgreater0. Because p\textgreater=7 all leading
  coefficients are units. Fact 2 forces 4e=3f; otherwise there is a
  unique least-weight term. Consequently v\_p(a) is even, say 2t, and
  v\_p(c\textsuperscript{2-b}2)=3t. The two differences involving a are
  units.
\item
  If p divides exactly two radii, say b,c, primitivity makes a a unit.
  By fact 3, unequal squared valuations make P nonzero, so
  v\_p(b)=v\_p(c)=t\textgreater0. After unit normalization write
  A=p\^{}\{2t\}U, B=p\^{}\{2t\}V. If U=V modulo p the degree-four part
  equals J U\^{}4 modulo p, which is nonzero; all terms of higher degree
  have strictly larger valuation. Consequently the two scaled squares
  differ modulo p, and v\_p(c\textsuperscript{2-b}2)=2t. The other
  differences are units.
\end{itemize}

No prime divides all three radii. The three gcds d\_a,d\_b,d\_c are
pairwise coprime; a shared prime has equal exponents in its two radii.
All remaining prime powers of each radius have even exponents. This
proves the displayed decomposition of the radii. By the primewise
analysis, primes dividing abc never divide g. For odd primes not
dividing abc the difference valuations are equal and exactly absorbed by
g. At 2 exactly one difference has one extra factor 2. For a prime
dividing a alone, its residual power in c\textsuperscript{2-b}2 is the
corresponding power of u\_a\^{}3; for a prime dividing b,c it is the
corresponding power of d\_a\^{}2. All other differences are units in
both cases. Permute the labels as needed. Unique factorization proves
the three identities. QED.

\subsection{Theorem 2 (pairwise coprime
radii)}\label{theorem-2-pairwise-coprime-radii}

If gcd(a,b)=gcd(a,c)=gcd(b,c)=1 after primitive integer normalization,
the original D=2 matrix is nonsingular.

Proof. If singular, Theorem 1 gives a=u\textsuperscript{2,b=v}2,c=w\^{}2
with positive u\textless v\textless w. The additive identity among
differences gives

\begin{verbatim}
epsilon_b v^3 = epsilon_a u^3 + epsilon_c w^3.
\end{verbatim}

Since w\textgreater v, epsilon\_b must be 2; the other epsilons are 1.
Thus v\^{}3 is the midpoint of u\^{}3 and w\^{}3. The function t
-\textgreater{} t\^{}(4/3) is strictly convex for t\textgreater0, so
w\textsuperscript{4-v}4\textgreater v\textsuperscript{4-u}4. On the
other hand the original difference equations imply
w\textsuperscript{4-v}4=g u\^{}3\textless g
w\textsuperscript{3=v}4-u\^{}4. Contradiction. QED.

\subsection{Corollary (only the two smallest radii share prime
factors)}\label{corollary-only-the-two-smallest-radii-share-prime-factors}

Suppose gcd(a,c)=gcd(b,c)=1 but d\_c=gcd(a,b)\textgreater1. The original
matrix is also nonsingular. Otherwise a=d\_c u\^{}2,b=d\_c
v\textsuperscript{2,c=w}2, hence usqrt(d\_c)v. The necessary additive
equation would give epsilon\_b v\^{}3=epsilon\_a u\^{}3+epsilon\_c
d\_c\^{}2 w\^{}3. Its right side exceeds
d\_c\textsuperscript{(7/2)v}3\textgreater2v\^{}3, impossible. QED.

\subsection{Computations and limits}\label{computations-and-limits}

New companion certificate reconstructs both numerator polynomials,
checks all local support identities, checks both exact rational scalars
for 6106 pairwise coprime triples among 16648 primitive triples up to
height 50, and screens all 24754 primitive triples with radii coprime to
30 up to height 200 against Theorem 1. No triple survives the necessary
identities. Combined with the prior small-prime theorems this supplies a
finite theorem for all primitive triples up to height 200, conditional
on the stated proof dependencies; no new direct determinant checks were
made for the entire sample. The earlier 2-adic and small-prime
certificates were rerun, including their original residual five direct
determinant checks. No conclusion for all triples sharing primes in the
other patterns. Not a proof of unrestricted rational regularity. No new
literature search, full old suite replay, independent audit, or Windows
validation.

\newpage

\section{Supporting source 4:
symmetric\_seven\_single\_gcd\_proof.md}\label{supporting-source-4-symmetric_seven_single_gcd_proof.md}

\section{Elimination of single shared-gcd patterns in the symmetric
seven-node D=2
problem}\label{elimination-of-single-shared-gcd-patterns-in-the-symmetric-seven-node-d2-problem}

2026-11-20 JST --- PRIVATE / HOLD NOVELTY INCONCLUSIVE / REFERENCES
PARTIAL / NOT HUMAN PEER REVIEWED / All rights reserved.

\subsection{Dependencies and status}\label{dependencies-and-status}

Uses Theorem 1 of \texttt{symmetric\_seven\_global\_proof.md} in the
2026-11-20 handoff. The 52 manifest entries and the frozen v0.7 archive
SHA256 ef40748e44e13280e00b4db496a3296119fa1adb2659afb869e5865f3765c6f4
matched. The global, 2-adic, and small-prime certificates were replayed
successfully in this session. No full v0.7 suite replay, independent
audit or Windows run. No change to the frozen archive. These proposed
proofs depend on the necessary residue reduction and radius formulas,
not their converse.

Write d=d\_a=gcd(b,c), e=d\_b=gcd(a,c), f=d\_c=gcd(a,b), and
u=u\_a,v=u\_b,w=u\_c.~Assume singularity, so a=ef u\^{}2, b=df v\^{}2,
c=de w\^{}2, c\textsuperscript{2-b}2=g eps\_a d\^{}2 u\^{}3,
c\textsuperscript{2-a}2=g eps\_b e\^{}2 v\^{}3,
b\textsuperscript{2-a}2=g eps\_c f\^{}2 w\^{}3. Exactly one epsilon is
2; all others 1. All radii are coprime to 30, hence each nontrivial gcd
is \textgreater=7. Also a\textless b\textless c.~Divide the additive
relation by the positive common factor (def)\^{}(-3/2) to obtain eps\_b
e\textsuperscript{(7/2)b}(3/2)=eps\_a
d\textsuperscript{(7/2)a}(3/2)+eps\_c f\textsuperscript{(7/2)c}(3/2).

\subsection{Proposition 1: if e=1, then f=1 is necessary for
singularity}\label{proposition-1-if-e1-then-f1-is-necessary-for-singularity}

From the last equality, 2 b\^{}(3/2)\textgreater{}
f\textsuperscript{(7/2)c}(3/2), so f\^{}(7/2)\textless2. As f is a
positive integer coprime to 30, f=1. This also excludes
d\textgreater1,f\textgreater1,e=1, extending the earlier
smallest-pair-only result.

\subsection{Proposition 2: d=f=1, e\textgreater1 is
impossible}\label{proposition-2-df1-e1-is-impossible}

Then a=e u\textsuperscript{2,b=v}2,c=e w\^{}2, so u\textless w. The
difference equations give c\textsuperscript{2-b}2 \textless=
2(b\textsuperscript{2-a}2), hence c\textsuperscript{2\textless3b}2. But
eps\_b e\^{}2 v\^{}3=eps\_a u\^{}3+eps\_c
w\textsuperscript{3\textless4w}3. Thus
(b/c)\textsuperscript{3=v}6/(e\^{}3
w\textsuperscript{6)\textless16/e}7\textless=16/7\^{}7\textless1/27, so
b/c\textless1/3, contradicting c\textsuperscript{2\textless3b}2.

\subsection{Proposition 3: d\textgreater1, e=f=1 is
impossible}\label{proposition-3-d1-ef1-is-impossible}

Here a=u\^{}2,b=d v\^{}2,c=d w\^{}2, with w\textgreater v. The additive
equality forces eps\_b=2, eps\_a=eps\_c=1. Then d\^{}2
u\textsuperscript{3=2v}3-w\textsuperscript{3\textless v}3. Put
X=(a/b)\^{}2, t=w/v, Z=(c/b)\textsuperscript{2=t}4. In particular
0\textless X\textless d\^{}(-14/3)\textless1/2401. Dividing the first
difference by the third produces
(t\textsuperscript{4-1)/(1-X)=(2-t}3)/t\^{}3,
t\textsuperscript{7=2-X(2-t}3). As 1\textless t\^{}3\textless2,
2-X\textless t\^{}7\textless2. Since (11/10)\^{}7\textless2-1/2401 and
(11/10)\^{}4\textgreater7/5, we obtain 7/5\textless Z\textless2.

The two scalar numerator polynomials P1,P2 in squared radii are
homogeneous, so a singular point would satisfy P\_i(X,1,Z)=0 for at
least one i. Reconstruct from the radius formulas. For each i the signed
numerator is p0(Z)+X p1(Z)+X\^{}2 p2(Z)+R(X,Z), where p0=-K
Z\textsuperscript{4(Z-1)}4; p1=L Z\textsuperscript{3(Z-1)}4(Z+1);
p2=-Z\textsuperscript{2(Z-1)}4(H Z\^{}2+I Z+H). (K,L,H,I)=(5,17,16,53)
or (3,11,12,37). All remaining powers of X are at least 3. On the
rectangle 0\textless X\textless=1/2401, 7/5\textless=Z\textless=2,
p1/\textbar p0\textbar\textless8 and p2\textless=0. Direct exact
expansion yields a bound \textbar R\textbar\textless=M X\^{}3, with
M=sum\_\{i\textgreater=3,j\}\textbar c\_ij\textbar2\^{}j=64184 or 44392.
In particular P \textless=
-K(7/5)\textsuperscript{4(2/5)}4(1-8/2401)+M/2401\^{}3 \textless{} 0.
Both final rational inequalities and the coefficient identities are
checked by the new companion certificate. Thus neither numerator can
vanish, contradiction. No finite radius bound is used in the proof.

\subsection{Proposition 4: an additional two-shared-gcd
subcase}\label{proposition-4-an-additional-two-shared-gcd-subcase}

If d=1, e\textgreater1, f\textgreater1 and eps\_a=2, no singularity is
possible. Here v\^{}2\textgreater e u\^{}2 because a\textless b. Define
S=(c\textsuperscript{2-b}2)/(c\textsuperscript{2-a}2)=2u\textsuperscript{3/(e}2v\^{}3).
Then 0\textless S\textless2/e\^{}(7/2). Also
b\textsuperscript{2/c}2=1-S(1-a\textsuperscript{2/c}2)\textgreater1-S.
The additive equality with eps\_b=eps\_c=1 gives
f\textsuperscript{2w}3=e\textsuperscript{2v}3(1-S). Consequently
(f/e)\textsuperscript{7=(b/c)}3(1-S)\textsuperscript{2\textgreater(1-S)}(7/2)\textgreater=1-(7/2)S\textgreater1-1/e.
The middle inequality is Bernoulli's inequality and the last uses
e\textgreater=7. But (f/e)\^{}7\textless1, so f\textless e and
(f/e)\textsuperscript{7\textless=(1-1/e)}7\textless=1-1/e,
contradiction.

\subsection{Remaining cases}\label{remaining-cases}

The earlier pairwise-coprime proof and Propositions 1--3 exclude all
patterns with at most one nontrivial shared gcd, and the pattern
d\textgreater1,f\textgreater1,e=1. Proposition 4 excludes part of
d=1,e\textgreater1,f\textgreater1. Any singular triple must have
e\textgreater1 and at least one of d,f\textgreater1. The remaining cases
are d=1,e\textgreater1,f\textgreater1 with eps\_a=1;
d\textgreater1,e\textgreater1,f=1; and d,e,f\textgreater1. The complete
rational seven-node theorem remains open in this work, as do the
unrestricted rational problem and the novelty assessment. The new
certificate checks exact polynomial identities and inequalities, not an
independent audit. No claim of first discovery.

\newpage

\section{Publication integrity and
limitations}\label{publication-integrity-and-limitations}

This document is an assembled proof-source edition; the previously
completed independent AI audit was performed on the source chain, not on
this PDF's typeset text. The original parent ZIP and Stage 05 audit
request ZIP have now been recovered; their SHA256 digests match the
audit records. The completion proof and certificate are byte-identical
across the recovered originals and the compiled source bundle. The four
major proof sources in this edition match the parent ZIP byte-for-byte.
The optional single-GCD source is supplementary historical work, not
needed for the final theorem. A verified frozen v0.7 original is
supplied as an immutable attachment.

The author has not yet authorized public GitHub/Zenodo release of this
prepared candidate. DOI not assigned. All rights reserved.

\end{document}
